% Source-faithful massless form and explicit unproved source-specific fields.
\chapter[Original Energy and Fermionic Growth]{Unnormalized Original Energy and the Fermionic Growth Problem}
\label{ch:energy}

\section{Why an operator construction does not yet imply confinement}
Chapter~\ref{ch:hamiltonian} established a positive self-adjoint physical
Hamiltonian from the \emph{original} closed sixteen-charge column,
\(H_{\mathrm{phys}}=\frac1{16}T^\dagger T\). The spectral question is more
demanding. We must understand how its quadratic form behaves as the bosonic
configuration escapes through the noncompact space
\(\Bconfiguration=\mathbb R^{9\times3}\). The scalar potential is quartic,
but it has noncompact valleys along commuting matrix configurations. On those
valleys it cannot alone force the first radial moment
\(\int\|X\|\,\|\Psi(X)\|^2\,dX\) to remain finite or bounded on a
form-norm ball. The original fermionic field is unbounded and linear in \(X\),
so dropping it is not legitimate.

The historical research took two converging paths. In October, FCP proved the
exact \emph{source-defined} undeformed charge/core-form identities (G4--K9D)
and constructed the actual closed physical operator and its domain (BA--BH).
In November, we returned to the energy decomposition and proved an explicit
original-source smooth-core composition, then built first-moment compactness
machinery without presupposing the missing physical confinement estimate.
These paths meet at a precise outstanding mathematical interface, not at the
mere similarity of their formulas.

\section{The source charge form and its massless energy identity}
The starting point is the original \texttt{AlgebraData N}, not an independent
choice of gamma matrices or a convenient replacement Fock space. Write
\(Q_\alpha^M\) for the literal smooth-core source supercharges from
Chapter~\ref{ch:charges} and
\(\iota:\mathscr D\to\OriginalLtwo\) for the source compact-core inclusion.
For \(f\in\mathscr D\), define
\begin{equation}\label{eq:energy-coreForm}
 q_M^{\mathrm{core}}(f):=\frac1{16}\sum_{\alpha=0}^{15}
   \|\iota(Q_\alpha^M f)\|_{L^2}^2.
\end{equation}
This is the original Lean declaration \texttt{M.coreForm f}; the coefficient
\(1/16\) is an exact normalization, not a proportionality convention.

\begin{theorem}[Original massless core-energy identity; integration
strict-qualified]\label{thm:massless-energy}
For every admissible source \(M:\texttt{AlgebraData N}\) and every original
\(f:\texttt{SmoothCore N}\), the actual source definitions satisfy
\begin{align}\label{eq:massless-energy}
 q_M^{\mathrm{core}}(f)
  &=\frac12\sum_{(i,a)\in\{0,\ldots,8\}\times\mathrm{Color}(N)}
     \|\iota(\partial_{i,a}f)\|_{L^2}^2 \\
  &\quad+\int_{\Bconfiguration_N}\!\left(
     V_M^{(1,0)}(X)\,\|f(X)\|^2
     +\operatorname{Re}\langle f(X),B_M^{(1,0)}(X)f(X)\rangle_{\mathbb C}
     \right)\,dX .\nonumber
\end{align}
Here \(V_M^{(1,0)}\) and \(B_M^{(1,0)}\) denote, \emph{without changing
any coefficients}, \texttt{M.deformedBosonicPotential 1 0} and
\texttt{M.deformedFermionField 1 0}. The displayed real part records
Lean's real inner product of the complex Fock fiber.
\end{theorem}
\begin{proof}[Accepted source-rewrite argument]
The October G4--K9D source theorem establishes, on every
original compactly supported smooth function, equality between
\texttt{M.deformedCoreEnergy 1 0} and \texttt{M.coreForm}. It uses the
original ordered gamma/color/fermion finite sums, including the massless
\(1/2\) coefficient; no added potential-agreement premise remains.
The existing source theorem \texttt{M.averaged\_core\_energy 1 0 f}
expands that \emph{same} deformed core energy into kinetic, bosonic and
fermionic components. Rewriting the left-hand side by G4--K9D then
applying the source expansion yields \eqref{eq:massless-energy}.
The integrated Lean proof is exactly this two-step composition; it does not
expand a second, independently stipulated supercharge.
\end{proof}
\formalrecord{\texttt{OAI.BFSSQuantum.AlgebraData.massless\_coreForm\_energy},
\texttt{SpectralMasslessCoreFormIdentity.lean}, blob
\texttt{3a8dd27bb573e62c8f86b79b83e0ab96f391602a}; strict
\href{https://github.com/etblink/math/actions/runs/42325087575}{\#42325087575}
(success). The historical FCP ancestor G4--K9D was independently accepted at
\href{https://github.com/etblink/Foundational-Convergence-Program/actions/runs/38011995220}{\#38011995220}.
Atlas A-01; original-source definition D-03.}

\begin{remark}[Scope of the identity]
The theorem is generic in \(M\). It has not, \emph{by itself},
(a) constructed a new Spin(9) action on the physical Hilbert space,
(b) identified every concrete Pauli-color contraction with a separately
printed wedge formula, (c) bounded the original field in operator norm,
or (d) transferred the formula from the smooth core to a high-Spin reducing
sector of the \emph{closed} physical Hamiltonian.
These are logically separate statements. The BA--BH theorems guarantee the
correct closed Hamiltonian exists, but do not make an arbitrary core-level
inequality invariant under closure and projection.
\end{remark}

\section{The quartic bosonic potential and its valleys}
For the usual concrete normalized Pauli \(\mathrm{SU}(2)\) color data,
write a bosonic configuration as three row vectors
\(X^a\in\mathbb R^9\), \(a=1,2,3\). The familiar wedge expression is
\begin{equation}\label{eq:potential-wedges}
 V_{\wedge}(X)=\sum_{1\le a<b\le3}\|X^a\wedge X^b\|^2
 =\sum_{a<b}\bigl(\|X^a\|^2\|X^b\|^2-\langle X^a,X^b\rangle^2\bigr).
\end{equation}
The source audit locates the concrete equality
\(V_M^{(1,0)}=V_{\wedge}\) in a normalized-Pauli color theorem whose
\emph{hypothesis must be bound to the actual paired instance}. We use
\eqref{eq:potential-wedges} here as the explicitly specified concrete
potential, not as a new blanket consequence of
Theorem~\ref{thm:massless-energy} for every \(M\).

The geometric obstruction is immediate. Choose a unit spatial vector
\(e\in\mathbb R^9\), a unit color vector \(n\in\mathbb R^3\), and a
radius \(r\ge0\). The rank-one configurations
\begin{equation}\label{eq:flat-valley}
  X^a=r n_a e,\qquad a=1,2,3,
\end{equation}
satisfy \(X^a\wedge X^b=0\) for every pair, hence
\(V_{\wedge}(X)=0\) for arbitrarily large \(r\). The quartic potential
is positive away from the rank-one locus, but not coercive in \(\|X\|\)
on its own. Rank-one \emph{configurations} are not normalizable physical
states; nonetheless they identify escape channels that any rigorous
spectral argument must control.

\section{The unnormalized fermionic field is a different object}
The original source fermionic field \(B_M^{(1,0)}(X)\) is an
endomorphism of the genuine \(2^{24}\)-dimensional Fock fiber in the
paired \(N=2\) specialization. Its coefficient construction is linear in
\(X\), using the original real gamma matrices, oriented normalized
\(\mathrm{SU}(2)\) structure constants and self-adjoint Majoranas with
\(\{\theta_p,\theta_q\}=\delta_{pq}I\). An often useful skew-CAR
representation is schematically
\begin{equation}\label{eq:B-skew-CAR}
 B(X)=\frac{i}{2}\sum_{p,q}K_{pq}(X)\theta_p\theta_q,
 \qquad K(X)^\mathsf T=-K(X), \quad K(tX)=tK(X).
\end{equation}
The exact source-index orientation and coefficient agreement between
\(B_M^{(1,0)}\) and \eqref{eq:B-skew-CAR} are an explicit \emph{binding
obligation}; a matching numerical norm is insufficient to establish the
signed operator identity.

If a source-typed bound
\begin{equation}\label{eq:B-linear-target}
 \|B_M^{(1,0)}(X)\|_{\mathrm{op}}\le C\|X\|
 \quad(X\in\mathbb R^{27}),\qquad 0\le C<\infty
\end{equation}
were established for the actual concrete data, its quadratic-form
contribution would satisfy
\begin{equation}\label{eq:B-negative}
 \operatorname{Re}\langle v,B_M^{(1,0)}(X)v\rangle
 \ge -C\|X\|\,\|v\|^2.
\end{equation}
The proof is the elementary estimate
\(|\langle v,Bv\rangle|\le\|B\|\|v\|^2\).
Thus a \emph{bounded} multiplier is not required here: a linear-growth
multiplier is what the original energy argument needs.

\begin{proposition}[Conservative finite-CAR growth estimate;
analytical and not Lean-qualified]\label{prop:B-analytic}
Under the exact 48-generator CAR, signed source-field identification,
and the skew-matrix trace identity recorded in the November 6 A/B audit,
one obtains the conditional conservative estimate
\begin{equation}\label{eq:analytic-C}
   \|B(X)\|_{\mathrm{op}}\le 8\sqrt3\,\|X\|.
\end{equation}
This is a conventional finite-dimensional derivation contingent on the
specified identification, not a newly accepted Lean source theorem.
\end{proposition}
\begin{proof}[Structure of the audited calculation]
Real orthogonal skew normal form decomposes the \(48\times48\) coefficient
matrix into 24 two-dimensional rotation blocks with nonnegative
frequencies \(a_1,\ldots,a_{24}\). CAR normalization
\(\theta_p^2=I/2\) makes the magnitude of the contribution from each
commuting fermion block \(a_j/2\). Hence
\begin{equation}\label{eq:B-HS-bound}
  \|B(X)\|_{\mathrm{op}}
  \le \frac12\sum_{j=1}^{24}a_j
  \le\frac12\sqrt{24}\left(\sum_j a_j^2\right)^{1/2}.
\end{equation}
The A/B source-normalization analysis gives
\(\sum_j a_j^2=32\|X\|^2\), which yields
\(\frac12\sqrt{24\cdot32}=8\sqrt3\).
The equality between the coefficient matrix and the \emph{original}
source field, and the source-level operator inequality, remain to be
formally checked. These qualification requirements are not supplied by
repeating the finite-dimensional arithmetic.
\end{proof}
\formalrecord{Analytical finite-source reconstruction:
\path{BFSS_SPECTRAL_CONFINEMENT_AB_AUDIT_2026_11_06.md} and
\path{BFSS_SPECTRAL_SOURCE_PREMISES_AUDIT_2026_11_06.md};
source-typed growth obligation O-02 in the atlas. Their numerical checks
are diagnostic evidence, not Lean acceptance.}

\section{Adversarial valley packets: a conditional noncompactness test}
\label{sec:valley-packets}
A rank-one commuting configuration is not itself a quantum state. A
serious test for a putative \emph{unsectored} coercive estimate is a
family of normalized, gauge-physical vectors $\psi_R$ whose mass moves
to distance $R\to\infty$ along the commuting valley and for which
\begin{equation}\label{eq:escape-form-test}
  \sup_R q_{\mathrm{phys}}(\psi_R)<\infty,\qquad
  \int\|X\|\|\psi_R(X)\|^2\,dX\longrightarrow\infty.
\end{equation}
If such vectors belonged to the original form domain and a claimed
coercivity inequality held uniformly on their subspace, the two
conclusions in \eqref{eq:escape-form-test} would contradict it.
The heuristic construction uses a slowly varying longitudinal packet
and the transverse supersymmetric oscillator ground band, whose
linear zero-point may cancel the fermionic term. \emph{This is not a
constructed original physical-domain sequence}: gauge averaging,
normalization, valley-band dependence, actual operator domains and
remainder bounds must be proved. Moreover, \emph{bounded quadratic-form
energy is not Weyl's criterion}; a genuine Weyl sequence for an
operator at energy $E$ additionally requires
$\|(H_{\mathrm{phys}}-E)\psi_R\|_2\to0$ (along with
normalization and weak escape). Neither claim is made here.
A verified low-type escaping sequence need not contradict a
high-Spin-type reducing-sector estimate; the exact isotype of the
constructed vectors is decisive.

\section{Why the first radial moment is the critical target}
Combining Theorem~\ref{thm:massless-energy} with a \emph{hypothetically
proved} source bound \eqref{eq:B-linear-target} would give, schematically,
\begin{equation}\label{eq:weak-valley-bound}
 q_M^{\mathrm{core}}(f)\ge
 \frac12\|\nabla f\|_2^2 +\int V_M^{(1,0)}(X)\|f(X)\|^2\,dX
   -C\int\|X\|\|f(X)\|^2\,dX.
\end{equation}
The negative final term can dominate along the flat valley, so
\eqref{eq:weak-valley-bound} does \emph{not} prove confinement. The
prospective high-Spin mechanism must rule out sufficiently many low
transverse oscillator levels on the \emph{whole} appropriate isotypic
component, thereby adding a positive linear-in-radius term that can
outweigh the fermionic growth. This is the task for
Chapter~\ref{ch:spectral}; its group/weight/sector assumptions remain
explicitly conditional.

There is a practical prioritization lesson. For qualitative confinement,
a rigorously source-typed \emph{finite} constant \(C\) in
\eqref{eq:B-linear-target} may suffice if a genuine family of high-type
physical sectors with sufficiently strong oscillator exclusion is
available. Obtaining the sharp coefficient \(8\sqrt3\) first may
therefore be unnecessary. But changing \(C\) changes the required
oscillator threshold, and that threshold cannot be adjusted without
checking actual source branching, nonemptiness and center conventions.

\begin{openproblem}[Physical unnormalized field and energy-to-moment
bridge]\label{op:energy-to-moment}
On the literal paired \(\mathrm{SU}(2)\) physical Hilbert space:
(i) prove the exact Pauli/gamma/ordered-Majorana binding of the source
fields, (ii) prove a source-typed operator growth bound, (iii) construct
the physically correct Spin(9) reducing action and high-weight exclusion,
and (iv) derive from the \emph{same} closed physical form a uniform
positive first-radial-moment and gradient estimate. Theorem~\ref{thm:massless-energy}
alone establishes none of (i)--(iv).
\end{openproblem}

\editornote{Evidence statuses: the source core form A-01 is integrated
strict-qualified; the concrete conservative norm constant is conditional
finite-dimensional analysis; the source physical growth and coercivity
joins O-02 and O-04 remain open. The normalized Gram observable of
Chapter~\ref{ch:norm} is auxiliary and does not replace the unnormalized
fermionic field.}
